\SourceLeaf{52}
\[
\begin{array}{lrl}
\log(n-1)\quad .\ .\ .\ .&0.0505666&\\
\log\cos v\quad .\ .\ .\ .\ .&9.6627075&\\
\log\operatorname{tang}h.\quad .\ .\ .\ .&0.0113399&\\
\log\operatorname{cotang}(F-\delta)\quad .\ .&0.0646895&\\\cline{1-2}
\log\cos w\quad .\ .\ .\ .\ .&9.7893035,&\text{whence }w=+52^0\ 0'\ 14''25\\
&&\text{and }u=10^0\ 36'\ 44''54
\end{array}
\]
\[
\begin{array}{lrl}
\log\operatorname{tang}h.\quad .\ .\ .\ .&0.0113399&\\
\log\cos u\quad .\ .\ .\ .\ .&9.9925074&\\\cline{1-2}
\log\operatorname{tang}G\quad .\ .\ .\ .&0.0188325,&\text{whence }G=46^0\ 14'\ 30''77\\
&&\text{and }G-\delta=37^0\ 51'\ 47''67
\end{array}
\]
\[
\begin{array}{lrl}
\log\cos G\quad .\ .\ .\ .\ .&9.8398647&\\
\log\operatorname{tang}u.\quad .\ .\ .\ .&9.2726964&\\
\text{Comp. }\log\sin(G-\delta).&0.2119881&\\\cline{1-2}
\log\operatorname{tang}\lambda.\quad .\ .\ .\ .&9.3245492,&\text{whence }\lambda=11^0\ 55'\ 18''31
\end{array}
\]
\[
\begin{array}{lrl}
\log\cos\lambda\quad .\ .\ .\ .\ .&9.9905300&\\
\log\operatorname{cotang}(G-\delta)\quad .\ .&0.1093281&\\\cline{1-2}
\log\operatorname{tang}\varphi\quad .\ .\ .\ .&0.0998581,&\text{whence }\varphi=51^0\ 31'\ 47''19
\end{array}
\]

This elevation of the pole differs from the true one by only a few seconds of arc. From $\lambda$ we further find $\gamma=307^0\ 17'\ 24''91$, or $20^{\mathrm h}\ 29^{\mathrm m}\ 9^{\mathrm s}66$ in time. The resulting lead of the clock over true sidereal time is $=+10^{\mathrm m}\ 58^{\mathrm s}34$, which agrees to within half a second with the determination made from other observations on the same day.
\Article{7}
Our investigation would certainly be incomplete unless we also show how large an error in the determination of the unknown quantities can result from the small, unavoidable errors made in measuring the altitudes. This will also show what kind of stars should preferably be chosen so that the determinations are as accurate as possible. For this purpose let us regard the altitudes $h,h'$ and the quantities $\varphi$ and $\lambda$ as variable. Differentiation of equation $1$ then gives
\expandafter\def\csname HistoricalNoteText1\endcsname{We always assume that azimuths are counted from the south to the right through west, north, and so on, from $0$ to $360^0$.}
\SourceLeaf{53}
\[
\cos h\,\mathrm dh=(\sin\delta\cos\varphi-\cos\delta\sin\varphi\cos\lambda)\mathrm d\varphi-\cos\delta\cos\varphi\sin\lambda\,\mathrm d\lambda
\]
Now it is known that if the azimuth of the vertical circle in which the first star was observed is set $=A$, we have\HistoricalNote{1}
\[
\begin{aligned}
\sin\delta\cos\varphi-\cos\delta\sin\varphi\cos\lambda&=-\cos h\cos A\\
\cos\delta\sin\lambda&=\cos h\sin A
\end{aligned}
\]
so the preceding equation becomes
\[
{}[17]\qquad\mathrm dh=-\cos A\,\mathrm d\varphi-\cos\varphi\sin A\,\mathrm d\lambda
\]
In exactly the same way, if the azimuth of the second star in the second observation is set $=A'$, differentiation of equation $2$ gives
\[
{}[18]\qquad\mathrm dh'=-\cos A'\,\mathrm d\varphi-\cos\varphi\sin A'\,\mathrm d\lambda
\]
Combining these equations gives the following equations.
\[
{}[19]\qquad\mathrm d\varphi=-\frac{\sin A'}{\sin(A'-A)}\mathrm dh+\frac{\sin A\;{}'}{\sin(A'-A)}\mathrm dh'
\]
\[
{}[20]\qquad\cos\varphi\,\mathrm d\lambda=\frac{\cos A'}{\sin(A'-A)}\mathrm dh-\frac{\cos A}{\sin(A'-A)}\mathrm dh'
\]

It follows at once that even the slightest errors in the observed altitudes can lead to quite considerable errors in determining both the elevation of the pole and the time if $\sin(A'-A)$ is very small, that is, if the second vertical circle is very close to the first or to the circle opposite it. We must therefore take care to avoid pairing two stars in such a position. It is also clear that whenever $\sin(A'-A)$ has a considerable positive or negative value, the observational errors can certainly increase only slightly in the determination of the elevation of the pole. The same holds for the determination of the time, whose precision also depends on the elevation of the pole itself. This arises from the nature of the problem itself, irrespective of the method, since the time is always determined with less precision at higher latitudes than in places closer to the equator. — The best choice would always be to select %
\SourceLeaf{54}stars such that $A'-A$ or $A-A'$ is approximately $=90^0$, that is, such that the two vertical circles intersect approximately at right angles. In this case it is easy to show that the maximum value of $\mathrm d\varphi$ is $=\BareRoot{(\mathrm dh^2+\mathrm dh'^2)}$ and that for the other maximum we have $\mathrm d\lambda=\frac{\BareRoot{(\mathrm dh^2+\mathrm dh'^2)}}{\cos\varphi}$. Thus if we could be sure of each altitude to within $10''$, the maximum error in determining latitude would be $=14''$ and the maximum error in determining the angle $\lambda$ would be $\frac{14''}{\cos\varphi}$. In our regions the maximum error in determining the time would therefore be $=1''5$ of time. It is hardly necessary to point out that in this estimate of the errors we consider only the effect of errors made in the observations. This is quite sufficient, since in our day the positions of the brighter stars themselves can be given with such precision that their errors can be disregarded.

\medskip\centerline{\rule{.25\textwidth}{.3pt}}\medskip
In the lectures I shall give in the coming winter semester

\noindent\textit{publicly}, at hour IX, I shall explain the theory of the parabolic motion of comets.

\noindent\textit{privately}, at hour X, I shall teach \textit{Astronomy}, also giving those attending the opportunity to do practical exercises both in observation and in astronomical calculation.
\par\medskip\centerline{\rule{.3\textwidth}{.3pt}}
